\documentclass[10pt,a4paper]{amsart}
\usepackage{amsmath,amssymb,amsthm,graphicx}
\usepackage{verbatim}
\newenvironment{reference}{\comment}{\endcomment}
%\newenvironment{reference}{}{}
\newenvironment{slogan}{\comment}{\endcomment}
\newenvironment{history}{\comment}{\endcomment}

% For commutative diagrams we use Xy-pic
\usepackage[all]{xy}

% We use 2cell for 2-commutative diagrams.
\xyoption{2cell}
% Theorem environments.
%
\theoremstyle{plain}
\newtheorem{theorem}[subsection]{Theorem}
\newtheorem{proposition}[subsection]{Proposition}
\newtheorem{lemma}[subsection]{Lemma}

\theoremstyle{definition}
\newtheorem{definition}[subsection]{Definition}
\newtheorem{example}[subsection]{Example}
\newtheorem{exercise}[subsection]{Exercise}
\newtheorem{situation}[subsection]{Situation}

\theoremstyle{remark}
\newtheorem{remark}[subsection]{Remark}
\newtheorem{remarks}[subsection]{Remarks}

\numberwithin{equation}{subsection}

% Macros
%
\def\lim{\mathop{\mathrm{lim}}\nolimits}
\def\colim{\mathop{\mathrm{colim}}\nolimits}
\def\Spec{\mathop{\mathrm{Spec}}}
\def\Hom{\mathop{\mathrm{Hom}}\nolimits}
\def\Ext{\mathop{\mathrm{Ext}}\nolimits}
\def\SheafHom{\mathop{\mathcal{H}\!\mathit{om}}\nolimits}
\def\SheafExt{\mathop{\mathcal{E}\!\mathit{xt}}\nolimits}
\def\Sch{\mathit{Sch}}
\def\Mor{\mathop{\mathrm{Mor}}\nolimits}
\def\Ob{\mathop{\mathrm{Ob}}\nolimits}
\def\Sh{\mathop{\mathit{Sh}}\nolimits}
\def\NL{\mathop{N\!L}\nolimits}
\def\CH{\mathop{\mathrm{CH}}\nolimits}
\def\proetale{{pro\text{-}\acute{e}tale}}
\def\etale{{\acute{e}tale}}
\def\QCoh{\mathit{QCoh}}
\def\Ker{\mathop{\mathrm{Ker}}}
\def\Im{\mathop{\mathrm{Im}}}
\def\Coker{\mathop{\mathrm{Coker}}}
\def\Coim{\mathop{\mathrm{Coim}}}

% Boxtimes
%
\DeclareMathSymbol{\boxtimes}{\mathbin}{AMSa}{"02}

%
% Macros for moduli stacks/spaces
%
\def\QCohstack{\mathcal{QC}\!\mathit{oh}}
\def\Cohstack{\mathcal{C}\!\mathit{oh}}
\def\Spacesstack{\mathcal{S}\!\mathit{paces}}
\def\Quotfunctor{\mathrm{Quot}}
\def\Hilbfunctor{\mathrm{Hilb}}
\def\Curvesstack{\mathcal{C}\!\mathit{urves}}
\def\Polarizedstack{\mathcal{P}\!\mathit{olarized}}
\def\Complexesstack{\mathcal{C}\!\mathit{omplexes}}
% \Pic is the operator that assigns to X its picard group, usage \Pic(X)
% \Picardstack_{X/B} denotes the Picard stack of X over B
% \Picardfunctor_{X/B} denotes the Picard functor of X over B
\def\Pic{\mathop{\mathrm{Pic}}\nolimits}
\def\Picardstack{\mathcal{P}\!\mathit{ic}}
\def\Picardfunctor{\mathrm{Pic}}
\def\Deformationcategory{\mathcal{D}\!\mathit{ef}}
\usepackage{cleveref}
\crefname{originalreference}{original reference}{original references}
\Crefname{originalreference}{Original reference}{Original references}
\makeatletter
\newlabel{definition-etale}{{00U1}{source}{}{sourcepointer-1}{}}
\newlabel{definition-etale@cref}{{[originalreference][0][]00U1}{[1][0][]source}{}{}{}}
\makeatother
\crefname{originalreference}{original reference}{original references}
\Crefname{originalreference}{Original reference}{Original references}
\makeatletter
\newlabel{lemma-functoriality-etale-site}{{03G2}{source}{}{sourcepointer-2}{}}
\newlabel{lemma-functoriality-etale-site@cref}{{[originalreference][0][]03G2}{[1][0][]source}{}{}{}}
\makeatother
\crefname{originallemma}{original lemma}{original lemmas}
\Crefname{originallemma}{Original lemma}{Original lemmas}
\makeatletter
\newlabel{sites-lemma-cocontinuous-morphism-topoi}{{00XO}{source}{}{sourcepointer-3}{}}
\newlabel{sites-lemma-cocontinuous-morphism-topoi@cref}{{[originallemma][0][]00XO}{[1][0][]source}{}{}{}}
\makeatother
\crefname{originallemma}{original lemma}{original lemmas}
\Crefname{originallemma}{Original lemma}{Original lemmas}
\makeatletter
\newlabel{sites-lemma-composition-cocontinuous}{{03L5}{source}{}{sourcepointer-4}{}}
\newlabel{sites-lemma-composition-cocontinuous@cref}{{[originallemma][0][]03L5}{[1][0][]source}{}{}{}}
\makeatother
\crefname{originallemma}{original lemma}{original lemmas}
\Crefname{originallemma}{Original lemma}{Original lemmas}
\makeatletter
\newlabel{sites-lemma-localize-given-products}{{03CE}{source}{}{sourcepointer-5}{}}
\newlabel{sites-lemma-localize-given-products@cref}{{[originallemma][0][]03CE}{[1][0][]source}{}{}{}}
\makeatother
\crefname{originallemma}{original lemma}{original lemmas}
\Crefname{originallemma}{Original lemma}{Original lemmas}
\makeatletter
\newlabel{sites-lemma-when-shriek}{{00XR}{source}{}{sourcepointer-6}{}}
\newlabel{sites-lemma-when-shriek@cref}{{[originallemma][0][]00XR}{[1][0][]source}{}{}{}}
\makeatother
\crefname{originalsection}{original section}{original sections}
\Crefname{originalsection}{Original section}{Original sections}
\makeatletter
\newlabel{sites-section-localize}{{00XZ}{source}{}{sourcepointer-7}{}}
\newlabel{sites-section-localize@cref}{{[originalsection][0][]00XZ}{[1][0][]source}{}{}{}}
\makeatother
\title{An etale morphism of algebraic spaces is given by a continuous and cocontinuous functor on the small etale sites}
\author{The Stacks Project Authors}
\date{stacks-project:spaces-properties.tex:local-snapshot-20261012; GNU FDL 1.2 or later}
\begin{document}
\maketitle
\noindent\textit{Editorial note.} All delivered mathematical text is the author's own, copied from the pinned Stacks Project chapter file spaces-properties.tex; no statement or proof step was rewritten, reordered or removed.\par\medskip
\noindent\textit{Editorial note.} The counted proof is the author's complete argument: the verification that the rule on objects defines a continuous and cocontinuous functor, the identification of its extension with a localization functor and the resulting morphism of topoi, and the comparison with the small morphism through the composition lemma together with the pullback formula.\par\medskip
\noindent\textit{Editorial note.} The label lemma-etale-morphism-topoi names the counted statement, and the proof environment immediately following it in the source is delivered verbatim as its proof.\par\medskip
\noindent\textit{Editorial note.} The counted statement is carried by the author's own lemma environment in the chapter source.\par\medskip
\noindent\textit{Editorial note.} Pre-open gate: python operations/agent-authoring/tools/precard\_gate.py --source spaces-properties.tex --label lemma-etale-morphism-topoi --start 3021 --end 3071 --local-id 28054 returned RESULT=OK (L28054 free, occupancy FREE, public interval clear). Candidate taken from sources/stacks-spaces-properties/candidates-halfA.json.\par\medskip
\noindent\textit{Editorial note.} The statement and proof are copied from the pinned chapter file; the cited inputs are kept as external pointers with their canonical Stacks tags.\par\medskip
\noindent\textit{Allowed standard prerequisites (editorial).}\begin{itemize}
\item Definition of an etale morphism of algebraic spaces. Same chapter, the definition of an etale morphism of algebraic spaces through the representability of the base change and etaleness of the induced morphisms of schemes. Supplies the fact that the rule j is well defined on the category of objects etale over the source.
\item Localization functor and the morphism of topoi it induces. Sites, the section on localization, and the lemma that a localization functor and the projection induce the same morphism of topoi. Supplies the identification of j extended to the big site with a localization functor and the resulting morphism of topoi.
\item Composition of cocontinuous functors and associated morphisms of topoi. Sites, the lemma that a commutative square of cocontinuous functors gives a commutative diagram of the associated morphisms of topoi. Supplies the commutativity of the diagram of topoi used to identify the morphism induced by j with the small morphism.
\item Construction of the small morphism of topoi. Same chapter, the lemma constructing the functoriality of the small etale site and the associated morphism of topoi. Supplies the description of the morphism f\_\{small\} with which the morphism induced by j is compared.
\end{itemize}\medskip
\ifdefined\hypertarget\hypertarget{sourcepointer-1}{}\fi
\noindent\textit{Source reference.} Original reference 00U1: sources/stacks-spaces-properties/pinned-original-source.tex, source lines 2312--2312. Prerequisite pointer invoked at proof line(s) 3040; it names a standard fact of the same chapter that the counted argument uses, and every step of the counted proof built on it is retained verbatim in the delivered body.\par\medskip
\ifdefined\hypertarget\hypertarget{sourcepointer-2}{}\fi
\noindent\textit{Source reference.} Original reference 03G2: sources/stacks-spaces-properties/pinned-original-source.tex, source lines 2826--2826. Prerequisite pointer invoked at proof line(s) 3070; it names a standard fact of the same chapter that the counted argument uses, and every step of the counted proof built on it is retained verbatim in the delivered body.\par\medskip
\ifdefined\hypertarget\hypertarget{sourcepointer-3}{}\fi
\noindent\textit{Source reference.} Original lemma 00XO: sources/stacks-chow/tags.txt, source lines 1226--1226. Prerequisite pointer invoked at proof line(s) 3031; it names a standard fact of a sibling chapter that the counted argument uses, and every step of the counted proof built on it is retained verbatim in the delivered body.\par\medskip
\ifdefined\hypertarget\hypertarget{sourcepointer-4}{}\fi
\noindent\textit{Source reference.} Original lemma 03L5: sources/stacks-chow/tags.txt, source lines 4701--4701. Prerequisite pointer invoked at proof line(s) 3066; it names a standard fact of a sibling chapter that the counted argument uses, and every step of the counted proof built on it is retained verbatim in the delivered body.\par\medskip
\ifdefined\hypertarget\hypertarget{sourcepointer-5}{}\fi
\noindent\textit{Source reference.} Original lemma 03CE: sources/stacks-chow/tags.txt, source lines 4381--4381. Prerequisite pointer invoked at proof line(s) 3063; it names a standard fact of a sibling chapter that the counted argument uses, and every step of the counted proof built on it is retained verbatim in the delivered body.\par\medskip
\ifdefined\hypertarget\hypertarget{sourcepointer-6}{}\fi
\noindent\textit{Source reference.} Original lemma 00XR: sources/stacks-chow/tags.txt, source lines 1229--1229. Prerequisite pointer invoked at proof line(s) 3033; it names a standard fact of a sibling chapter that the counted argument uses, and every step of the counted proof built on it is retained verbatim in the delivered body.\par\medskip
\ifdefined\hypertarget\hypertarget{sourcepointer-7}{}\fi
\noindent\textit{Source reference.} Original section 00XZ: sources/stacks-chow/tags.txt, source lines 1237--1237. Prerequisite pointer invoked at proof line(s) 3061; it names a standard fact of a sibling chapter that the counted argument uses, and every step of the counted proof built on it is retained verbatim in the delivered body.\par\medskip
\begin{lemma}
\label{lemma-etale-morphism-topoi}
Let $S$ be a scheme, and let $f : X \to Y$ be a morphism of algebraic spaces
over $S$. Assume $f$ is \'etale. In this case there is a functor
$$
j : X_\etale \to Y_\etale, \quad
(\varphi : U \to X) \mapsto (f \circ \varphi : U \to Y)
$$
which is cocontinuous. The morphism of topoi $f_{small}$ is the
morphism of topoi associated to $j$, see
Sites, Lemma \ref{sites-lemma-cocontinuous-morphism-topoi}.
Moreover, $j$ is continuous as well, hence
Sites, Lemma \ref{sites-lemma-when-shriek}
applies. In particular $f_{small}^{-1}\mathcal{G}(U) = \mathcal{G}(jU)$
for all sheaves $\mathcal{G}$ on $Y_\etale$.
\end{lemma}

\begin{proof}
Note that by our very definition of an \'etale morphism of algebraic spaces
(Definition \ref{definition-etale}) it is
indeed the case that the rule given defines a functor $j$ as indicated.
It is clear that $j$ is cocontinuous and continuous, simply because a covering
$\{U_i \to U\}$ of $j(\varphi : U \to X)$ in $Y_\etale$ is the
same thing as a covering of $(\varphi : U \to X)$ in $X_\etale$. It
remains to show that $j$ induces the same morphism of topoi as $f_{small}$.
To see this we consider the diagram
$$
\xymatrix{
X_\etale \ar[r] \ar[d]^j &
X_{spaces, \etale} \ar@/_/[d]_{j_{spaces}} \\
Y_\etale \ar[r] &
Y_{spaces, \etale} \ar@/_/[u]_{v : V \mapsto X \times_Y V}
}
$$
of categories. Here the functor $j_{spaces}$ is the obvious extension of $j$
to the category $X_{spaces, \etale}$. Thus the inner square is
commutative. In fact $j_{spaces}$ can be identified with the
localization functor
$j_X : Y_{spaces, \etale}/X \to Y_{spaces, \etale}$
discussed in
Sites, Section \ref{sites-section-localize}.
Hence, by
Sites, Lemma \ref{sites-lemma-localize-given-products}
the cocontinuous functor $j_{spaces}$ and the functor $v$ of the diagram
induce the same morphism of topoi. By
Sites, Lemma \ref{sites-lemma-composition-cocontinuous}
the commutativity of the inner square (consisting of cocontinuous functors
between sites) gives a commutative diagram of associated morphisms of topoi.
Hence, by the construction of $f_{small}$ in
Lemma \ref{lemma-functoriality-etale-site} we win.
\end{proof}

\end{document}
